\section{Discussion}
\label{sec:discussion}
\subsection{Separating auxiliary supervision from forecast readout}
Across the four selected seeds, auxiliary future-state supervision increased the direct-readout mean Macro AUROC, AP, and F1, and explicit mean forecast readout increased those means further. These averages preserve unfavorable individual results: auxiliary-supervised direct AUROC was below \nojepa{} at seed 2026, and forecast-mean AUROC was below direct readout at seed 42. A favorable metric average does not imply improvement in every training realization or at every operating point.

The full-versus-\nojepa{} Macro AUROC comparison met both the nine-test principal-family criterion (0.00450) and the all-78 sensitivity criterion (0.03850). Its AP and F1 differences did not meet the principal-family criterion, and no full-versus-direct Macro comparison did so. In the uniform 99,999-resample table, T1 and T2 AUROC comparisons against \nojepa{} also passed secondary-family and all-78 adjustment. However, an independent 999,999-resample check changed the T2 all-78 value from 0.04408 to 0.054872 when the other 77 values were held fixed. The T2 global decision is therefore sensitive to Monte Carlo precision; the check did not repeat the Macro or T1 simulations. These findings support conditional AUROC differences among the selected fitted predictors; they do not independently confirm a general architectural benefit or remove retrospective model and seed selection bias.

The \direct{}--\nojepa{} contrast most closely isolates future-state supervision. The full-versus-direct contrast jointly changes access to predicted state and change features and classifier input width. Classifier-width-matched direct readout, stopped-gradient forecast features, copy-state readout, and independent holdout forecast evaluation would help distinguish these contributions. The main model uses equal spatial weights in its forecast readout while retaining self-attention in the shared Transformers. It predicts embeddings rather than generating MRI for classification. The internal name \nojepa{} denotes removal of future-state supervision while retaining the shared fusion network and transition queries.

\subsection{What additional visits change}
Forecast mean's mean AUROC and F1 peaked at T1, whereas mean AP increased across T0--T2. Its paired T2--T0 intervals included zero for all three metrics. Changes in true- and false-positive predictions explain why more observations, ranking, and a fixed-threshold decision rule need not improve together. Outcome-conditioned histograms summarize probability changes within each seed-specific ensemble; they do not demonstrate that every patient's prediction improves or that score changes are treatment effects.

T0 imaging preceded randomization, while the model is conditioned on the eventual assigned arm. This remains a retrospective conditional scenario. The arm encoding does not capture delivered dose, cycle timing, dose reductions, adherence, or treatment changes. Altering the arm token cannot identify an individual treatment effect without suitable causal assumptions and exposure data. Complete T0--T2 holdout eligibility also restricts the population to which the longitudinal comparison applies.

\subsection{Classification, probability quality, and hypothetical utility}
The additional metrics expose trade-offs behind the ranking means. Compared with auxiliary-supervised direct readout, forecast mean had higher sensitivity, NPV, balanced accuracy, and MCC but lower specificity and slightly lower PPV and accuracy. Neither the sensitivity nor specificity contrast passed secondary-family adjustment. Their point estimates and unadjusted intervals remain relevant descriptions of the selected operating point, rather than multiplicity-adjusted evidence for a clinical benefit or harm. Both direct and forecast-mean accuracy were only marginally above the all-negative reference.

Forecast mean had worse NLL, Brier score, and ECE point estimates than both direct-readout controls. Its higher ranking means therefore did not imply better predicted probabilities. The OOF Youden analysis slightly increased F1 for \nojepa{} but decreased it for the other models. Development OOF scores originate from one held-out fold model per patient, whereas test scores average five models within a seed. Transferring thresholds between these score distributions is a sensitivity analysis rather than independent ensemble calibration.

At the illustrative thresholds 0.30, 0.50, and 0.70, direct readout had higher mean net benefit than forecast mean despite lower AUROC. These are comparisons under a generic positive-action utility, not evaluations of a defined care pathway. A clinically motivated threshold requires consequences for false-positive and false-negative actions to be specified and calibration to be fitted without holdout leakage.

\subsection{Seed variation, subgroups, and forecasting diagnostics}
The mean and sample standard deviation across four seeds describe four training realizations using identical development folds and holdout patients. They do not turn seeds into independent clinical replications. Patient-bootstrap intervals and tests answer a different question: how the mean metric of these fixed predictors varies under patient resampling. Because seeds are not resampled and models are not refitted, those intervals do not estimate uncertainty over the training algorithm or future seed choices.

Subgroup summaries are exploratory, use small denominators, and do not test subtype interactions. The HR-positive/HER2-negative and HR-negative/HER2-negative full-versus-control Macro AUROC intervals excluded zero without multiplicity correction, whereas the other two HR/HER2 intervals included zero. These findings do not establish a corrected subtype-specific benefit. Low sensitivity at 0.5 in the HR-positive/HER2-negative group also persists. Source namespaces share one clinical release and do not provide external validation. The crop subgroup with 11 patients and only one positive outcome cannot support reliable discrimination assessment under the minimum-outcome rule.

Archived fixed-training-probe diagnostics are separate from the pCR holdout analysis. Their copy-normalized skills and gradient norms are calculated in each snapshot's own EMA representation, averaged over folds within seed, and then summarized across seeds. They do not ensemble latent states or establish anatomical forecast fidelity. At epoch 20, auxiliary-supervised direct models had higher mean copy-normalized skill than forecast mean at all three horizons, despite the latter's higher mean pCR AUROC. This weakens an explanation based solely on more accurate forecasts. Gradient routing supports an implementation-level mechanism, but gradient magnitude is not a biological effect size. Independent forecast evaluation at the selected pCR checkpoints remains necessary.

\subsection{Limitations and a confirmatory study}
\label{sec:limitations}
The choice of forecast mean as the overall model and the seed set 42, 44, 71, and 2026 were made after earlier holdout results had been available. This model and reporting-policy change, probability ensembling, and the statistical procedures are retrospective. Patient-bootstrap inference conditions on the fitted ensembles and fixed OOF thresholds. It excludes retraining, overall-model and seed selection, aggregation choice, checkpoint selection, and threshold-selection uncertainty. Four selected seeds do not establish broad training stability, and their sample standard deviations have limited precision. The reported mean describes four five-model predictors; it is not the performance of one twenty-model ensemble or a model trained on the full development cohort.

The 95\% patient intervals are unadjusted. Holm corrections for nine principal and 69 secondary tests, together with the all-78 sensitivity analysis, do not undo retrospective choices. Null-centered bootstrap tests are approximate, with finite-sample limitations particularly for non-smooth statistics such as AP, F1, and binned ECE. The conditional AUROC finding and the absence of adjusted AP/F1 or operating-point differences should remain visible together. Nonsignificance does not establish equivalence. A confirmatory study should lock the overall model, seeds, aggregation, checkpoint and threshold selection, and statistical comparisons before evaluating independent patients.

The holdout is an internal downstream split. Twenty-nine patients had belonged to an earlier development split. Current supervised models were retrained using the new patient-disjoint folds and were not warm-started from those earlier pCR models, but codec and localization-model pretraining have not been fully audited for current-holdout overlap. This prevents a claim that every pipeline stage was evaluated on wholly unseen patients.

Fixed baseline localization and cropping may omit multifocal disease or changes outside the field of view. Geometry and array checks do not establish anatomical registration accuracy. Holdout patients have complete T0--T2 imaging, whereas development retains incomplete histories; missingness, acquisition variation, and subtype prevalence may limit transportability. Archived training probes do not replace best-checkpoint holdout forecast evaluation and have small long-horizon denominators.

The study remains an internal architectural comparison. Matched clinical-only, conventional imaging, and independent longitudinal baselines are still needed, along with external validation, anatomical quality review, and independently assessed calibration. Precise pCR pathology criteria, scan intervals, ethics determination, and upstream provenance also require final documentation before submission.
